% !TEX TS-program = xelatex
% 替换方括号中的内容。如需使用自己的照片，请上传 photo.jpg。
\documentclass{resume}
\defaultfontfeatures[FandolSong-Regular]{Extension=.otf,SmallCapsFont=FandolHei-Regular}
\usepackage[scheme=plain,fontset=fandol]{ctex}
\usepackage{tabularx}
\hypersetup{hidelinks}

\begin{document}
\pagenumbering{gobble}
\renewcommand{\arraystretch}{1.2}

\noindent
\begin{minipage}[c]{0.82\textwidth}
    \section*{\faGraduationCap\ 个人信息}
    \begin{tabularx}{\linewidth}{@{}lX lX@{}}
        姓名 & [姓名] & 院校 & [院校名称] \\
        邮箱 & [邮箱] & 电话 & [电话号码] \\
        % 主页 & \href{https://example.com}{[个人主页]} & & \\
    \end{tabularx}
\end{minipage}\hfill
\begin{minipage}[c]{0.12\textwidth}
    \resumephoto[照片]{photo.jpg}
\end{minipage}\par

\section{\faHome\ 教育背景}
\datedsubsection{\textbf{[院校名称]}}{[城市]}
\datedsubsection{\textit{[学历与专业]}}{[起止时间]}
\textbf{绩点：} [绩点] \qquad \textbf{排名：} [排名] \\
\textbf{相关课程：} [课程名称]

\section{\faHeartO\ 荣誉奖励}
\datedline{[奖项或奖学金名称]}{[年份]}

\section{\faUsers\ 科研与项目经历}
\datedline{\textbf{[项目名称]}}{[起止时间]}
\textbf{项目内容：} [简述项目背景与目标。] \\
\textbf{采用方法：} [概述主要方法或技术。] \\
\textbf{个人贡献：} [说明负责的工作。] \\
\textbf{项目成果：} [概括交付成果或可量化的成效。]

\section{\faCogs\ 技能与兴趣}
\textbf{技术技能：} [编程语言、工具或框架] \\
\textbf{语言能力：} [语种及熟练程度] \\
\textbf{研究兴趣：} [研究方向或职业兴趣]

\end{document}
